\documentclass[professionalfonts]{beamer}
\usepackage{cmap}
\usepackage{lmodern}
\usepackage{microtype}
\usepackage[english]{babel}
\usepackage{tikz}
\usepackage{minted}
\usepackage{xcolor}

\usetheme{AnnArbor}
\usecolortheme{spruce}
\setbeamercolor{button}{bg=black, fg=yellow}

\usetikzlibrary{positioning}
\usetikzlibrary{calc}

\title{Lustrean}
\subtitle{Lustre + Lean = Lustrean}
\author{Jean \textsc{Caspar} \and Adrien \textsc{Mathieu}}
\date{December 2024}

\makeatletter
\newcommand*{\currentname}{\@currentlabelname}
\def\ft{\frametitle{\currentname}}
\makeatother

\setbeamertemplate{itemize items}[square]
\setbeamertemplate{enumerate items}[default]

\begin{document}

\begingroup
\setbeamertemplate{headline}{}
\setbeamertemplate{footline}{}
\begin{frame}[noframenumbering]
  \titlepage
\end{frame}
\endgroup

\section{Hello}
\subsection{Bye}
\begin{frame}
  \ft{}
\end{frame}

\section{Abstract domains}
\subsection{Domain}
\begin{frame}
  \ft{}
  A domain over $n$ variables a data structure
  which represents an overapproximation of a
  set of environments. It has:

  \begin{itemize}
    \item decidable equality,
    \item a function which returns a new environment,
    \item an assignation function which takes an integer expression
    and updates the environment accordingly,
    \item a guard function which takes a boolean expression
    and refines the environment to keep only those states
    that makes the boolean exwhich should be different
    from $\bot$.pression be true,
    \item a widening and a narrowing operators.
  \end{itemize}

  The code is inspired from the TP of abstract interpretation
  from Antoine Miné, Marc Chevalier and Josselin Giet.
\end{frame}

\subsection{Non relational domain}
\begin{frame}
  \ft{}
  A non relational domain is a domain of the form
  $\{0, \dots, n-1\} \to V$ where $V$ is a domain
  representing values.

  The value domains are bounded lattice with a widening
  and a narrowing which support arithmetic operations.

  We did not use any relational domains, as it would have
  been to long to implement and to prove properties on them.
\end{frame}

\subsection{Nil in values domains}
\begin{frame}
  \ft{}
  Values domains are also required to contain an
  element representing $nil$. The order of the lattice
  is an order represeting the set of possible values
  the expression can take ; it is not the order
  ``x is less defined than y''. Especially, we do
  not have $nil \sqsubseteq 3$.
\end{frame}

\subsection{Concrete value domains}
\begin{frame}
  \ft{}

  \begin{itemize}
    \item Integers with a top and bottom element.
    \item Intervals.
  \end{itemize}
  These domains send nil to top.
  \begin{itemize}
    \item A domain combinator which transform a domain $D$
    into $D \times \mathbb{B}$, where $(v, \operatorname{false})$
    represents the set of values represented by $v$, except $nil$,
    and $(v, \operatorname{true})$ represents the set of values
    represented by $v$ and $nil$.
  \end{itemize}
\end{frame}

\section{Interpreter}
\subsection{CFG}
\begin{frame}
  \ft{}
  \begin{itemize}
    \item The interpreter is given a list of ``pre nodes'',
    parameterized by the number of variables.
    \item A pre node is a list of tuples $(inst, out, synt)$.
    \item $inst$ is an instruction, $out$ the index of the exit node
    as a natural number, $synt$ a pointer to the piece of code from
    which the instruction originates.
  \end{itemize}
\end{frame}

\begin{frame}
  \ft{}
  \begin{itemize}
    \item We compile the pre nodes to a cfg. It consists of an array of nodes and arcs.
    \item An arc contains an instruction, the index of the source node and of the exit node,
    and a pointer to the syntax.
    \item A node contains the list of the indices of the arcs going into it.
  \end{itemize}
  Indices lives in $Fin~arcs.size$, respectively $Fin~nodes.size$.

  Because we have to translate the first representation into the second while
  assuring correctness, there is a lot of bookkeeping
  involved with dependent types, and I used
  no less that two intermediate structures.

  We also compute a set of widening points, which select at least
  one node in each strongly connected component of the cfg.
\end{frame}

\subsection{Execution}
\begin{frame}
  \ft{}
  \begin{itemize}
    \item After that, we build a state which assigns a environment
    (living in a domain) to each node and arc.
    \item First we iterate over all arcs and assigns to each
    arc the value of the transfer function of this arc
    to the source environment.
    \item Then we iterate over the nodes and assigns
    to the them the union of the environments of their
    entrant arcs. If the node is a widening point, we perform a widening.
    \item We currently don't use the narrowing.
    \item When it stabilizes, we return the state.
  \end{itemize}
\end{frame}

\section{Correctness}
\begin{frame}
  \ft{}
  \begin{itemize}
    \item The notion of domain is totally formalized,
    save for the termination property of the narrowing
    and of the widening: we need classical logic to prove
    it, but then it would be difficult to extract the witness.
    \item We didn't provide a concrete semantics and proved
    that the abstract semantics is a sound abstraction of the concrete
    semantics.
    \item We don't have an other semantics to which we can
    compare our interpreter and relatively to which
    we can prove soundness.
  \end{itemize}
\end{frame}

\end{document}
